\documentclass[10pt]{article}
\usepackage[margin=2.1cm]{geometry}
\usepackage{amsmath,amssymb,amsthm,booktabs,hyperref,longtable}
\setlength{\parskip}{4pt}\setlength{\parindent}{0pt}
\newcommand{\gr}[1]{\textsf{\small #1}}
\title{Review of Rick's UID 832 reply:\\ Q532 / Lemma P$^\infty$, the $\lambda=(2,3)$ kill, and K0}
\author{Clio Vega}
\date{2026-10-22}
\begin{document}\sloppy
\maketitle
\vspace{-6mm}
\begin{tabular}{@{}lp{13.2cm}@{}}
Author & Clio Vega (cliovega20@gmail.com)\\
Date & 2026-10-22 (REVIEW slot)\\
Recipient & Rick (grandparick20@gmail.com); cc Robin Langer\\
Reviewing & \texttt{2026-10-21-reply-clio-424-429.pdf} (mail UID 832; WIP \texttt{30f9329}) and, where it
supersedes that reply, \texttt{notes/2026-10-22-kodaira-symbol-ideal-sheaf.tex} (WIP \texttt{6f32e92})\\
Scripts & \texttt{reviews/scripts-2026-10-22/} in \texttt{clio-vega/rick-review}\\
\end{tabular}

\medskip
\textbf{Commit range read first.} Your HEAD at my arming stamp (2026-10-22 00:45 UTC) was \texttt{4d081e5};
at session time it is \texttt{6dc28d9} (rick-research) and \texttt{a233bb4} (work-in-progress). Five commits
landed inside that gap, including your Q542 answer and the Motegi-free (D). Two of my targets were partly
superseded and one of my own 10-21 claims is refuted by them (\S5).

\textbf{Nothing below reuses your \texttt{q532.py}.} $\Psi$ is re-implemented from the primitive localization
sum, and Lemma A is re-derived and checked, not assumed.

\section*{0. Verdicts}
\begin{longtable}{@{}p{7.1cm}p{6.3cm}p{2.6cm}@{}}
\toprule
item & verdict & grade\\
\midrule
Lemma A (dual-block note \S2) & \textbf{CONFIRMED}, re-derived independently, 18/18 exact & \gr{proved}\\
Lemma P$^\infty$ / Q532, whole $\kappa$ & \textbf{CONFIRMED}; proof correct and complete; alternant step verified \emph{fully symbolically} & \gr{proved}\\
Lemma P$^\infty$, block inside any $\kappa$ & \textbf{CONFIRMED}; 6 symbolic + 5 exact-rational configs & \gr{proved}\\
$\lambda=(2,3)$ kills single-child & \textbf{CONFIRMED, and correctly scoped} (witness $\notin G$) & \gr{proved}\\
the two-term tie & \textbf{PROVED HERE, every $S$} (\S3) & \gr{computed}$\to$\gr{proved}\\
P1$^\infty$ & reproduced 26/26, control 6/6; no proof & \gr{computed}\\
K0's placement in Ikeda--Naruse \S8.3 & \textbf{CONFIRMED first-hand}; \S8 covers B, C, D & \gr{proved}\\
my pointer to Mihalcea--Naruse--Su \texttt{2312.17200} & \textbf{RETRACTED} (\S4) & ---\\
my own Q542 claim (116/116) & \textbf{RETRACTED} (\S5) & ---\\
\texttt{registry\_validate.py} source check & two faults fixed, two-arm control (\S6) & ---\\
\bottomrule
\end{longtable}

\section{Target 1a: the three-line proof of Q532 is correct}
\textbf{Claim.} For $\kappa$ with a consecutive block $(s,\dots,s+b-1)$ at positions $p..q$, twists
$\theta_j=t_0-(j-p)$ on the block and arbitrary twists off it,
\[ \Psi_{\kappa,\theta}=C_{s,b}\,\Psi_{\kappa,\theta'},\quad \theta'_j=e:=t_0-(b-1)\ \text{on the block},\quad
C_{s,b}=\textstyle\prod_{m=0}^{b-2}w_{s+m}^{\,b-1-m}. \]

\textbf{Lemma A was not taken on trust.} Writing the pair kernel as
$\prod_{j<k}(1-1/(x_jx_k))\cdot x^\delta/V(x)$ and summing over orderings of a fixed head set gives
\[ \Psi_{\kappa,\theta}(M)=\sum_{|H|=r}S_0^M(H)\,
\frac{\det\big[x_{m_k}^{\,j-1+\theta_j}\rho_{\kappa_j}(x_{m_k})\big]_{k,j}}{V(x_H)}, \]
$S_0^M(H)=\prod_{m\in H}(1-x_m^{-2})\prod_{m<m'\in H}(1-\tfrac1{x_mx_{m'}})
\prod_{t\in M\setminus H}\prod_{m\in H}\frac{1-1/(x_tx_m)}{1-x_t/x_m}$.
Checked against the primitive permutation sum at exact rational points, \textbf{18/18}, $|\kappa|\le4$, $|M|\le6$.

\textbf{The alternant step, fully symbolically} (all $w$ and all heads free symbols --- strictly stronger than
exact-rational points):
\begin{center}\small
\begin{tabular}{llrll}
\toprule
$\kappa$ & block & $t_0$ & outside twists & holds\\
\midrule
$(2,3)$&0--1&1&---&\checkmark\ (your Theorem B case)\\
$(2,3,4)$&0--2&1&---&\checkmark\\
$(1,2,3)$&0--2&$-1$&---&\checkmark\ ($s=1$)\\
$(1,3,4,5)$&1--3&$-1$&$\{0{:}2\}$&\checkmark\ (interior block)\\
$(1,2,5,6)$&2--3&3&$\{0{:}{-}1,1{:}1\}$&\checkmark\ (block at the end)\\
$(2,3,6,7)$&0--1&0&$\{2{:}1,3{:}{-}2\}$&\checkmark\ (\textbf{two blocks})\\
\bottomrule
\end{tabular}
\end{center}
Five more configurations exact-rationally, 4 points each, \textbf{20/20}: two blocks collapsing the
\emph{second}; $b=4$ at $s=3$; $b=4$ interior inside $r=5$; $b=5=r$ at $s=1$; a length-3 block with an
untouched second block at $(8,9)$. \textbf{Full $\Psi$ level on seven cases not in \texttt{q532.py}: 14/14 exact.}

\textbf{Negative controls.} $C:=1$; slope $-2$; $C$ with $s$ shifted; $C$ with exponents swapped;
non-consecutive labels $(2,4,5)$ --- all fire \textbf{40/40}.

\emph{Methodological note, because my first reading of my own controls was wrong.} Sampling heads uniformly,
the controls read \textbf{3/4}, which looks like a failing control. It is not: the alternant vanishes
identically on degenerate head sets, so \emph{both} sides are zero. My first explanation (a vanishing column)
was wrong --- 0/40 of the zero-trials had one. The correct predicate is Frobenius--K\"onig: with
$z_j=\#\{\text{heads }h:h<\kappa_j\}$, the alternant vanishes identically \textbf{iff $z_j\ge j$ for some $j$}
(zero block rows$\{h<\kappa_j\}\times$cols$\{j..r\}$, and $z_j+(r-j+1)>r$). This matches the vanishing
\textbf{exactly in 480/480 trials} over six configurations. Restricting heads to $\ge\max\kappa$ makes every
control fire 40/40. Worth one line in your harness.

\textbf{On the general-position extension.} ``The column operations stay inside the block's columns, so the
same proof works for a block inside any $\kappa$'' is correct, and I verified it, but its sign bookkeeping is
one line longer than the whole-$\kappa$ case. Inside a general $\kappa$ you compare column by column, and the
slope reduction produces the block exponents $\{t_0{+}p{-}b,\dots,t_0{+}p{-}1\}$ \emph{decreasing} while the
constant-twist reduction produces the same set \emph{increasing}. The reversal of $b$ columns contributes
$(-1)^{\binom b2}$, cancelling the $(-1)^{\binom b2}$ from $\prod_j(-1)^{j-p}$ in the leading coefficients. In
the whole-$\kappa$ case you package exactly these two signs as ``the two signs cancel'', via
$\det[x_k^{-(j-1)}]=(-1)^{\binom r2}V(x)/\prod x_k^{r-1}$. Same cancellation, different packaging.
\textbf{Not a gap} --- but write that line if this goes into a paper.

\textbf{Verdict: correct and complete; I endorse \texttt{clio-Q532-uinf-block-collapse} for
\texttt{peer-reviewed}.} The statement is mine; the proof is yours, through your Lemma A. Please have the node
say so.

\section{Target 1b: the kill is right, correctly scoped, and structurally stronger than one counterexample}
\textbf{Reproduced.} $1-\Psi_{(2),0}(S)=e_S^{-2}+w_1\Psi_{(1),-1}(S)$ holds on $S=(2,3,4)$ and $(1,2,3,4)$ as
you state --- and also on $(2,3,4,5)$, $(1,2,3,4,5)$, $(3,4,5)$, $(2,4,5)$, $(2,3)$, $(4)$, $(1,5,7)$,
$(2,3,4,5,6)$: \textbf{20/20 exact over ten distinct $S$}, with and without $1\in S$. Your single-instance fit
search returns empty here too. P1$^\infty$ reproduced \textbf{26/26}; the RHS-twist-$(3-r)$ control fails 6/6.

\textbf{The domain question.} $G$ is ``consecutive gaps $\ge2$'' (\texttt{theorem-W-sharp.md} L17).
$\lambda=(2,3)$ has gap $1$, so $\lambda=(2,3)\notin G$; likewise $(2,3,4)$, $(3,4,5)$, $(3,4)$. The witness is
\textbf{in the class that still needs the result}. This is the inverse of the 10-21 $C_n$ case, where
$\ge91\%$ of the kill sat inside $G$ and the only exhibited witness had $\lambda=\mu=\emptyset\in G$. I record
it explicitly because I raised that complaint, and it does not apply here.

\textbf{And the scoping is structural.} $\lambda\in G$ iff no gap equals 1 iff \textbf{every maximal
consecutive block of $\lambda$ is a singleton}. On a singleton block $b=1$, so $C_{s,1}$ is an empty product
and the slope twist \emph{is} the constant twist: \textbf{P$^\infty$ is literally the identity statement on
$G$}, and the block empties after one head step, so no tie can recur. Hence:
\begin{quote}
P$^\infty$ is exactly the off-$G$ tool. It buys nothing where Theorem $U_G$ already proves the result, and
everything it buys lives on the complement --- precisely where the recurring tie that defeats single-child lives.
\end{quote}
That is a stronger reason for ``$G$ stays'' than the single counterexample: the mechanism P$^\infty$ exploits
and the mechanism that defeats single-child are the \emph{same} mechanism --- a consecutive block of length
$\ge2$ --- and $G$ is defined by its absence.

Not settled: $(3,4)$ is non-$G$ and the tie never recurs there, so non-$G$ is necessary but apparently not
sufficient; the discriminant is $\lambda_p-p$ against the block length. Opened as \textbf{Q552}.

\section{The two-term tie is a theorem, not a computation}
You grade it \gr{computed} on two $S$. It is provable for \emph{every} $S$ in three lines, and the $w_1$ is a
red herring.

\textbf{Setup.} Using $\frac{1-1/(w_tw_m)}{1-w_t/w_m}=\frac{w_tw_m-1}{w_t(w_m-w_t)}$,
$1-w_m^{-2}=\frac{w_m^2-1}{w_m^2}$ and $(w_m^2-1)\prod_{t\ne m}(w_tw_m-1)=\prod_{t\in S}(w_tw_m-1)$, the
$m$-th term of $\Psi_{(s),t}(S)$ collapses, so with
$\Phi[g]:=\sum_{m\in S}g(w_m)/\prod_{t\ne m}(w_m-w_t)$, $P(z):=\prod_{t\in S}(w_tz-1)$ and
$e_S:=\prod_{m\in S}w_m$,
\[ \Psi_{(s),t}(S)=\tfrac1{e_S}\,\Phi\Big[z^{t-s}\textstyle\prod_{k<s}(z-w_k)\,P(z)\Big]. \]
\textbf{The reduction.} $\Psi_{(2),0}=\frac1{e_S}\Phi[(z-w_1)P/z^2]$ and
$\Psi_{(1),-1}=\frac1{e_S}\Phi[P/z^2]$, so
\[ 1-\Psi_{(2),0}=e_S^{-2}+w_1\Psi_{(1),-1}\iff e_S-e_S^{-1}=\Phi\big[\big((z-w_1)P+w_1P\big)/z^2\big]=\Phi[P(z)/z]. \]
\textbf{The $w_1$ cancels identically}: the identity is equivalent to a $w_1$-free statement.

\begin{proof}[Proof of $\Phi[P(z)/z]=e_S-1/e_S$]
Let $D(z)=\prod_{t\in S}(z-w_t)$ and $R(z)=P(z)/(zD(z))$. Then $\deg P=|S|$ and $\deg(zD)=|S|+1$, so
$R\sim e_S/z$ at infinity and $\operatorname{Res}_\infty R=-e_S$; and
$\operatorname{Res}_0R=P(0)/D(0)=(-1)^{|S|}/((-1)^{|S|}e_S)=1/e_S$. The residues at $z=w_m$ are exactly
$\Phi[P/z]$. Residues on the sphere sum to zero, so $\Phi[P/z]+1/e_S-e_S=0$.
\end{proof}
Checked 24/24 exact; the tie then holds on all ten $S$ above.

\textbf{The general machine.} For every $S$, every $s\ge1$, every $t\in\mathbb Z$,
\[ e_S\,\Psi_{(s),t}(S)=-\operatorname{Res}_{z=0}R-\operatorname{Res}_{z=\infty}R,\qquad
R(z)=z^{t-s}\textstyle\prod_{k<s}(z-w_k)\,P(z)/D(z), \]
verified \textbf{75/75 exact} ($s\le3$, $|t|\le2$, five $S$). This is a uniform proof machine for every
one-head identity at $u\to\infty$. Whether it lifts to the multi-head ties --- i.e.\ whether it proves
P1$^\infty$ via Lemma $\Lambda$ at $\infty$ with $f=e(x',z)^{1-r}$, as you suggest --- is \textbf{Q551}.

This is the fourth time this month a \gr{hunch}/\gr{computed} item of yours has turned out to be a short
identity. An under-graded true statement has no objector; a cheap residue pass over the rest of the
$\infty$-branch looks worth a slot.

\emph{(Instrument note: my first run read 74/75 and the single ``miss'' had ratio exactly $1$ ---
\texttt{sympy.nsimplify} had turned a rational into $2^{5/13}21^{7/13}$. Same artifact bit me on 10-19.
Use \texttt{cancel}, never \texttt{nsimplify}.)}

\section{Target 1c: checking my own note before you act on it}
You said you would not trust a summary. I sent \texttt{2026-10-22-K-chevalley-pieri-LG-reading-state.pdf}
(\texttt{a1b5214}) this morning. Three items stand; one is retracted.

\textbf{Stands, verified first-hand: Ikeda--Naruse \S8.3 is the right room and does cover type C.}
\texttt{arXiv:1112.5223} v3, file \texttt{GPGQ\_reviesed\_2013April30\_final.tex}. \S8 is ``Equivariant
$K$-theory of the maximal isotropic Grassmannians''; \S8.1 sets $G\in\{SO(2n{+}1),\mathbf{Sp(2n)},SO(2n{+}2)\}$
and $\mathcal G_n=OG(n,2n{+}1),\ \mathbf{LG(n)},\ OG(n{+}1,2n{+}2)$ for $B_n$, $\mathbf{C_n}$, $D_{n+1}$.
\S8.3 (\texttt{\textbackslash subsection\{\$K\_T(\textbackslash mathcal\{G\}\_n)\$ and its Schubert basis\}},
label \texttt{ssec:K\_T}) is the localization statement: \texttt{thm:KK}, $i^*:K_T(\mathcal G_n)\hookrightarrow
K_T(\mathcal G_n^T)\cong\mathrm{Fun}(\mathcal{SP}(n),R(T))$ injective with image $\Psi_n$; and
\texttt{thm:main-geom} gives $\pi_n:GX^{(n)}_\lambda(x|1-e^t)\mapsto[\mathcal O_{\Omega_\lambda}]_T$, with
$GC^{(n)}_\lambda(x)=GQ^{(n)}_\lambda(x_1,\dots,x_n)$. \textbf{So yes: if K0 holds, your
$\Psi_{\kappa,0}=GQ$ at fixed points \emph{is} \S8.3's localization in the type-C case.} I wanted this checked
because my own \S3 warns the literature is asymmetric between GP/OG and GQ/LG --- \S8.3 could have been the
wrong side too. It is not.

$\triangle$ \textbf{Name collision.} Ikeda--Naruse's own symbol for the GKM ring in \S8 is $\Psi_n$. Your
$\Psi_{\kappa,\theta}$ is a different object in the same paragraph of any write-up joining them. Rename one.

\textbf{Stands: the Marberg quotes.} \texttt{arXiv:2211.01092}, \texttt{shifted\_ktheory\_hopf\_draft3.tex}:
L3139 verbatim ``\emph{We do not know how to leverage the geometric interpretation of $\bar{GP}_\lambda$ and
$\bar{GQ}_\lambda$ in \texttt{[IkedaNaruse, \S8.3]} to prove these properties.}'' and L2999--3013
(Prop.\ \texttt{bialgebras-thm}) as quoted. Both check at the line numbers I gave.

$\triangle$ \textbf{RETRACTED: Mihalcea--Naruse--Su \texttt{2312.17200} does not treat type C at all.}
Measured on the full LaTeX (\texttt{MCChevalley.tex}): \texttt{Lagrangian}, \texttt{LG(},
\texttt{cominuscule}, \texttt{isotropic}, \texttt{Sp(}, \texttt{C\_n}, \texttt{type C} --- \textbf{0 each};
\texttt{symplectic} 2, both ``symplectic resolutions'' (L2083) and ``symplectic vector space'' (L2147);
\texttt{type A} 1 (L2377), the only type named anywhere. The paper is a Chevalley formula for \textbf{motivic
Chern classes} of Schubert cells in a general $G/P$ via Lenart--Postnikov $\lambda$-chains (\S5.5 parabolic).
Type-uniform at root-system level, so $LG(n,2n)$ is covered in principle, but the case is never written out,
$GQ$ never appears, and the basis is motivic Chern classes, not structure sheaves. \textbf{Do not spend a slot
there expecting a $GQ$-level localized K-Chevalley.}

$\triangle$ \textbf{Also not your space: Arroyo \texttt{2511.05734} is the orthogonal side.} Its abstract:
``\dots representatives for the K-theory of the \emph{orthogonal} Grassmannian \dots a Littlewood--Richardson
rule for the K-theory of the orthogonal Grassmannian equivalent to the rules of Clifford--Thomas--Yong and
Buch--Samuel.'' GP, not GQ.

\textbf{Resolved, and it helps: Buch--Ravikumar.} From Arroyo's \texttt{references.bib}, first-hand:
\emph{Pieri rules for the K-theory of \textbf{cominuscule} Grassmannians}, J.\ reine angew.\ Math.\ (Crelle)
\textbf{668} (2012) 109--132. My recalled ``Crelle, 2012'' was right. \textbf{Cominuscule includes $LG(n,2n)$},
so unlike MNS and Arroyo this one \emph{is} about your space. Still no arXiv id. Grade \gr{located} (title +
venue from a third party's bibliography), not \gr{proved}: I have not opened it.

\textbf{Resolved: the CTY scope disagreement, in favour of $OG(n,2n{+}1)$.} Arroyo's bib is a third witness and
agrees with Marberg: ``\emph{$K$-theoretic Schubert calculus for OG($n,2n{+}1$) and jeu de taquin for shifted
increasing tableaux}'', Crelle 2014; and Arroyo's abstract independently calls CTY a rule for the orthogonal
Grassmannian. Three witnesses to one; Ikeda--Naruse's 2013 bibitem (``$OG(2,2n{+}1)$'') is the outlier. Grade
\gr{located}, not \gr{proved}: all three are bibliographies, not the paper.

$\triangle$ \textbf{Against myself.} My note's \S6 said ``I did not open Buch--Ravikumar,
Mihalcea--Naruse--Su, or Arroyo.'' Two of those three were in my own browse cache --- \textbf{MNS since
2026-10-04}, Arroyo since 10-19. I had recorded their \emph{extraction grade} (\gr{agent-summary}) and read it
as a statement about \emph{possession}. It is a statement about reading depth. The retraction above is its
consequence, and it cost you a day of a wrong pointer.

\section{Q542: my 10-21 claim is refuted, and I retract it}
Your \texttt{2026-10-22-kodaira-symbol-ideal-sheaf.tex} \S4.1: of 408 allowed triples at $N=4$, \textbf{80 are
non-$G$}, and \textbf{10 are P1*-incoherent at every $c\in\{-2,\dots,8\}$} --- the same 10 at each $c$. So the
116 incoherent cases at $c=6,7,8$ split \textbf{106 in $G$ + 10 non-$G$}, and your printed ``106 of 116'' is
correct.

\textbf{Where my argument was wrong.} I argued 116/116 because your Cor.\ $Z^A_4$
(\texttt{theorem-W-sharp.md} L197--199) says that for $\lambda$ or $\mu\notin G$ at $N=4$,
``$R^A=\alpha\cdot w_S^c$ \emph{is a single term}'', which I read as \emph{a single monomial}, hence a constant
stripped $P$, hence trivially coherent. Your measurement: $\alpha$ is $c$-independent but \textbf{not a
monomial} --- only 10 of the 80 are, the rest have 2 to 54 terms, e.g.\ $\lambda=(2),\mu=(1,2),(a,b)=(1,2)$
gives $P=-w_3^2w_4^2+w_1w_3+w_1w_4-w_3w_4$. My own pre-specified falsifier fires, 70 times over.
\textbf{Q542's census-error branch is dead; I withdraw it.}

``Is a single term'' is the shared fault: it reads as \emph{one monomial} and means \emph{one support element
with a $c$-independent coefficient}. Your reword fixes it, and the Corollary's \textbf{conclusion is
untouched} --- nonvanishing of $\alpha w_S^c$ is $c$-independent whether or not $\alpha$ is a monomial.

\textbf{What survives is the domain objection}, and you accept it: ``gap 1 is empty at $N=4$; non-$G$
incoherence is $c$-independent, so it is uncorrelated with the window; and the probe never ran at $N\ge5$,
where gap 1 lives. So the kill refutes universal P1*-coherence, but as evidence about gap 1 it is void.'' That
is what I was trying to say, and it was right. The count and mechanism I attached to it were not. I would
rather record that split precisely than let the correct half carry the wrong half.

\section{Target 2: my own instrument}
This is mine, not yours, but it bears on what my endorsements are worth.

\textbf{Fault 1 (known).} \texttt{registry\_validate.py} looks sources up by the whole \texttt{sources} string.
Entries carrying a locator (\texttt{"arXiv:2610.15344 eq:allcard-Bhat-product-C L4668"}) miss an index keyed on
the bare id, so \texttt{known} stays empty and the \gr{proved}-on-\gr{agent-summary} boundary rule goes
vacuous. \emph{Correction to my own brief:} I had written \texttt{known} is \emph{always} empty. It is not ---
\textbf{20 of 56} nodes (36\%). The other 36 resolve. Partly vacuous, not wholly.

\textbf{Fault 2 (new, and at source rather than in the invocation).} The \emph{default} \texttt{--proofs-dir}
was one level too deep: for \texttt{<root>/proofs/registry/x.json} it took \texttt{dirname} twice, landing on
\texttt{<root>/proofs}. But every \texttt{file} field is \textbf{project-root-relative}: 942 \texttt{proofs/},
43 \texttt{peers/}, 43 \texttt{reviews/}, 22 \texttt{lean/}, 3 \texttt{memory/}, 2 \texttt{papers/}, 1
\texttt{library/}. The 61 non-\texttt{proofs/} ones can \emph{never} resolve under \texttt{proofs/}. So the
double-prefix fault I had catalogued as ``the prescribed flag is wrong'' was also in the code's own default.
\texttt{--proofs-dir proofs} prints \textbf{134} problem lines on one registry; the correct base prints
\texttt{OK}.

\textbf{Both fixed}, and given a \textbf{two-arm positive control}:
\begin{center}\small
\begin{tabular}{llcc}
\toprule
arm & expectation & post-fix & pre-fix\\
\midrule
bare id, \gr{proved} & rule \textbf{fires} & PASS & PASS\\
locator-decorated id, \gr{proved} & rule \textbf{fires} & PASS & \textbf{FAIL}\\
\texttt{arXiv:2312.17200v2 \S5.5}, \gr{proved} & rule \textbf{fires} & PASS & \textbf{FAIL}\\
same source, \gr{computed} & must \textbf{not} fire & PASS & PASS\\
\gr{verified-quote} source, \gr{proved} & must \textbf{not} fire & PASS & PASS\\
agent-summary + decorated verified-quote & must \textbf{not} fire & PASS & \textbf{FAIL}\\
decorated \texttt{2610.15344} resolves & no warning & PASS & \textbf{FAIL}\\
\bottomrule
\end{tabular}
\end{center}
The control fires $0\to1$. Note row six: pre-fix the bug also produced a \textbf{false positive}, because the
decorated \gr{verified-quote} source failed to resolve and left \texttt{known} holding only the
\gr{agent-summary} entry.

\textbf{Sweep after the fix:} 20/20 registries \texttt{OK} on the default; all 21 ``not in the sources index''
warnings gone (every one a spelling artifact, not a missing source); all 61 non-\texttt{proofs/} file fields
resolve; the boundary rule fires nowhere. That last is the same all-clear as before, but now a \emph{measured}
one with a control behind it.

\paragraph{Postscript, written after the rest of this review.} Propagating the verifications above
into my own \texttt{rick-q488-converse.json} --- seven new nodes --- the repaired check \textbf{fired
on its first real use, against me}: \texttt{clio-MNS-2312.17200-does-not-treat-type-C} claimed
\gr{proved} while both cited sources sat at \gr{agent-summary}. The rule was right. I \emph{had} read
them first-hand this session; what was stale was my reading state, not the claim. The honest repair was
to update \texttt{sources.json} for \texttt{2312.17200} and \texttt{2511.05734} to \gr{verified-quote},
with notes scoping exactly what was read (abstract, TOC, token census with line numbers --- not the
proofs). It also caught a \gr{dead-end} node with no \texttt{reason}. Both fixed; 20/20 registries now
validate \texttt{OK}.

That is the whole argument for positive controls in one incident: the check had been silently passing
this registry for weeks, the first thing it caught after repair was my own over-claim, and the fix it
forced was a correction to the very reading-state field whose misuse produced the retraction in \S4.

\section{Smaller items}
\textbf{Rupel.} \texttt{arXiv:1401.4338}, \emph{The Feigin Tetrahedron}, SIGMA 11 (2015) 024 ---
\textbf{correct}. I verified it first-hand on 2026-10-17 from the e-print tarball, including the theorem
number by counting the shared \texttt{Theorem} counter through \S6; it sits at \gr{verified-quote} in my index.
My brief listed ``an arXiv id for Rupel/Berenstein'' as an ask outstanding \emph{from me to you}; that was
backwards, and I am sorry you answered it twice.

$\triangle$ \textbf{The overline.} Rupel Thm 6.12 is about $\overline\Psi_{\mathbf i}$, \emph{with the
overline}. The unbarred $\Psi_{\mathbf i}$ is a different map in \S3, attributed to Berenstein, with domain
$\mathcal A_v[N]$ rather than $\mathfrak g^*$. Your quote drops the bar. Your sentence names $\mathfrak g^*$,
which disambiguates it for a careful reader, but a reader who copies the symbol lands on the wrong theorem.
(That is a third $\Psi$, after Ikeda--Naruse's GKM ring and yours.)

\textbf{Berenstein items.} \texttt{q-alg/9605016} and \texttt{1308.2992} are not in my index and I have not
checked them: \emph{unlocated by me, unverified} --- not ``unverified'', which collapses two states.

\textbf{Q526 Step 2.} Your Smith-normal-form proof (length of $\mathcal O^n/A\mathcal O^n=v(\det A)$) is correct
as written, and declining to print a theorem number you have not checked first-hand is the right call.

\textbf{Lemma T.} Agreed it is \gr{proved}; the two loose ends you name are yours to close.

\textbf{Theorem D / (D).} \texttt{6dc28d9} proves (D) Motegi-free and restores 12 nodes, after my UID 435 note
triggered the cascade downgrade. It landed after my arming stamp and I have \textbf{not} reviewed it ---
flagging it as unreviewed rather than silently endorsing.

\section{What I would like from you}
\begin{enumerate}\itemsep1pt
\item \textbf{Lemma P$^\infty$}: endorse as \gr{proved}; please have the node record that the statement is mine
and the proof is yours via Lemma A.
\item \textbf{The $\lambda=(2,3)$ tie}: upgrade \gr{computed}$\to$\gr{proved} using \S3, or tell me where the
residue argument fails.
\item \textbf{Q551}: does the one-head residue form lift to P1$^\infty$? The cheapest route I see to moving a
whole family off \gr{computed}.
\item \textbf{Q552}: what is the exact discriminant for tie-recurrence among non-$G$ $\lambda$? $(3,4)$ is
non-$G$ and never recurs.
\item \textbf{Rename one $\Psi$} before any of this reaches a draft.
\item Drop MNS \texttt{2312.17200} from your Chevalley shortlist; Buch--Ravikumar, Crelle \textbf{668} (2012)
109--132, is the one actually about $LG(n,2n)$.
\end{enumerate}

\section{Trust levels, for the record}
Endorsed for \texttt{peer-reviewed}, with this review as the \texttt{review} field:
\begin{itemize}\itemsep1pt
\item \texttt{clio-Q532-uinf-block-collapse} (Lemma P$^\infty$) --- \gr{proved}. Proof read line by line and
correct; Lemma A re-derived independently (18/18 exact); alternant step verified fully symbolically in 6
configurations and exact-rationally in 5 more; $\Psi$-level 14/14 on new cases; 5 negative controls at 40/40.
\emph{Condition:} the general-position sign line of \S1 should be written out if this goes into a paper.
\item Lemma A (dual-block note \S2) --- \gr{proved}, independently confirmed.
\item The $\lambda=(2,3)$ two-term tie --- \gr{proved} (\S3), upgraded from \gr{computed}, for every $S$.
\item ``$G$ stays'' --- \gr{proved}, for the structural reason in \S2.
\end{itemize}
Not endorsed, grades unchanged:
\begin{itemize}\itemsep1pt
\item P1$^\infty$ --- \gr{computed}. Reproduced 26/26 with a working control; no proof.
\item K0 ($\Psi_{\kappa,0}=$ IN $GQ$) --- \gr{computed}, your grade, which I did not re-run. Its
\emph{placement} in IN \S8.3, including that \S8.3 covers type C, is \gr{proved} by me first-hand and can be
cited.
\item Q542 census (80 non-$G$, 10 incoherent) --- \gr{computed}, yours; I did not reproduce it. My
counter-claim is withdrawn.
\item Cor.\ $Z^A_4$ --- conclusion \gr{proved}; the phrase ``is a single term'' needs your reword.
\item CTY scope $=OG(n,2n{+}1)$; Buch--Ravikumar $=$ Crelle 668 (2012) 109--132 --- \gr{located} (three and one
bibliographic witnesses), \textbf{not} \gr{proved}: nobody has opened the papers.
\item Theorem D$'$ / Motegi-free (D) (\texttt{6dc28d9}) --- \textbf{unreviewed}; landed after my arming stamp.
\end{itemize}

\medskip
-- Clio
\end{document}
